\section{Conclusão}

Este artigo apresentou uma infraestrutura reprodutível para descoberta de instâncias do SAPL, extração de PLs, tradução de ementas em ações textualizadas, busca semântica e associação exploratória com indicadores oficiais municipais. Na execução congelada de 12 de novembro de 2025, o sistema encontrou 1.259 instâncias do SAPL, extraiu PLs de 322 delas e reuniu um acervo de 220.065 proposições distribuídas por 322 municípios e 19 UFs. Os resultados demonstram viabilidade técnica para descoberta em escopo nacional, textualização em escala e utilidade prática da busca semântica quando a indexação usa ações textualizadas. No benchmark automatizado com 15 problemas e 272 julgamentos produzidos por um agente GPT-5.4, a representação em ações superou a busca sobre ementas em nDCG@10, MAP@10 e recuperação de documentos fortemente relevantes. A leitura mais segura desses resultados, porém, é de vantagem média no benchmark exploratório avaliado, e não de superioridade universal para qualquer domínio ou consulta. Mesmo sem substituir avaliação humana especializada, essa linha de base automatizada já oferece referência suficiente para comparar a qualidade relativa entre os dois métodos. O próximo passo é ampliar o benchmark com mais problemas e avaliadores, incluir \textit{baselines} lexicais, híbridos e de fusão, reduzir a endogeneidade do pool de avaliação e aumentar o controle sobre ruídos da base e das métricas observacionais.
